% Introduction for "The Unchosen Adjacency" (revised)

\section{Introduction}

Two ordinary events on a large social platform look unrelated. In the first, a user publishes a post or Reel, and advertisements the user did not select can be placed among the user's posts, over the Reel, or immediately after it, where, on the controls reviewed, an ordinary user cannot prevent them. In the second, a user blocks another person to stop a stream of hostile comments. The block succeeds, and it also does something the user did not ask for: it removes the other person's ability to see anything the user does, including the conduct the user meant to keep visible. Intermediate controls exist, and Section~\ref{sec:apps} examines what they do and do not separate.

The first event adds a relation the user did not choose. The second removes a relation the user meant to keep. Both are cases of overreach: a single action, offered through a coarse interface, changes several relations at once, and the user is left with the combination. A third failure runs the other way, underrealization: an action may fail to make a change the user did intend, as when a restriction is imposed without the explanation or end date that correction requires. The essay states this structure exactly, examines it on one large platform with another as comparator, and argues that a good deal of what users experience as platform overreach follows from it.

The thesis is that platforms govern partly by deciding which relations can be changed separately. Their power lies not only in permitting or forbidding actions but in packaging actions into bundles whose meaning the user does not control. A user may want to speak without advertising, exclude someone from a conversation without excluding them from view, or take part in a network without lending it approval. Where the interface offers no action that does exactly this, the interface decides what the user's act means.

The vocabulary matters. \emph{Compelled association} names the additive case, where an unwanted relation is attached. It also names a doctrine in constitutional law, which presupposes state action and does not apply to a private platform that users can leave. The essay does not rely on that doctrine. It uses \emph{unchosen bundling} as the general term for overreach, with compelled association and compelled dissociation as its two directions, and it treats the availability of exit as part of the phenomenon: exit is the only disaggregation on offer, and its cost is what the absence of finer controls imposes. The title names the additive case, the advertisement attached to a post, while the general term is bundle.

Two features distinguish the account from a complaint about bad design. First, it distinguishes mismatches that matter from mismatches that do not. Every software action has side effects and limits, so the account asks further questions: whether observers read a side effect as the user's own act, whether the user loses something they value or is made to bear a relation they reject, and whether the user fails to obtain a change they rate highly. Second, it is comparative. An interface is condemned only where a cleaner action is technically feasible, as shown by other systems, and is not offered.

\paragraph{Related work.} The nearest conceptual neighbor is contextual integrity, which asks whether information flows accord with the norms of their context \cite{nissenbaum2004}. The present account asks a different question, whether distinct relational changes can be selected separately, and it covers relations that are not information flows, such as an added advertisement or a removed view. Context collapse names the merging of several audiences into a single context \cite{marwick-boyd2011}. The collapse studied here is of the action space, in which several changes are tied to one act, and not of the audience. The graduated moderation sequence in Section~\ref{sec:moderation} parallels Ostrom's principle of graduated sanctions, under which sanctions are proportional to the severity of the violation \cite{ostrom1990}, and Grimmelmann's taxonomy of moderation techniques (exclusion, pricing, organizing, and norm-setting) \cite{grimmelmann2015} situates moderation as governance. The constitutional doctrine mentioned above belongs to the law of freedom of association, as in \emph{Roberts v.\ United States Jaycees} \cite{roberts1984}, and it constrains the state.

The argument proceeds as follows. Section~\ref{sec:defs} states the definitions: signed relational changes, surplus and deficit, attribution and value, and the notion of a defective action. Section~\ref{sec:apps} applies them to blocking, advertising adjacency, continued participation, and recommendation, and shows that the framework returns different provisional verdicts for different cases and excludes at least one. Section~\ref{sec:moderation} develops the constructive case, in which regulating participation is separated from regulating perception. Section~\ref{sec:objections} addresses objections. Section~\ref{sec:withdrawal} states a hypothesis about withdrawal, that users who cannot disaggregate expression from unwanted association will reduce the expression, and identifies what evidence would test it. Section~\ref{sec:design} lists candidate separable controls.

The essay makes no claim that any particular platform acts in bad faith or that the interface choices it discusses are accidental. The claim is structural: where meaning is bundled, someone other than the user is setting it, and the user's remaining recourse is to act less or to leave.
